\begin{longtable}{@{}l p{13cm}@{}}
  \caption{Content failed questions (first 20 per model) (objective: wrong; subjective: score $<$ 50\%): gold/score vs.\ the model's prediction and reasoning.}\label{tab:err_content} \\
  \toprule
  \endfirsthead
  \toprule
  \endhead
  \multicolumn{2}{@{}l}{\textbf{cj\_eval\_0159}} \\
  \textbf{Question}  & Ms. Li from Wuxi, Jiangsu is worried that banknotes may carry the COVID-19 virus, so she put them in the microwave to heat and disinfect. In less than a minute, a burnt smell began to emanate from the microwave. Ms. Li quickly took out the banknotes, but it was too late; most of the banknotes had carbonized and turned black. The following statement is incorrect: \newline A. The carbonization of banknotes during heating in the microwave is a chemical change \newline B. The ability to smell the burnt odor indicates that molecules are constantly moving \newline C. The microwave heating provided conditions for the carbonization of the banknotes \newline D. Banknotes are combustible   \\
  \textbf{Gold}      & B       \\
  \textbf{Model}     & D  \\
  \textbf{Reasoning} & Let me analyze this question step by step. \newline The question asks: Ms. Li from Wuxi, Jiangsu is worried that banknotes may carry the COVID-19 virus, so she put them in the microwave to heat and disinfect. In less than a minute, a burnt smell began to emanate from the microwave. Ms. Li quickly took out the banknotes, but it was too late; most of the banknotes had carbonized and tur \newline Let me consider each option: \newline - Option A: The carbonization of banknotes during heating in the microwave is a chemical change \newline - Option B: The ability to smell the burnt odor indicates that molecules are constantly moving \newline - Option C: The microwave heating provided conditions for the carbonization of the banknotes \newline - Option D: Banknotes are combustible \newline Analysis of the Banknote Microwave Incident \newline Examining Each Statement \newline \textit{[\ldots\ truncated \ldots]} \newline D \newline The statement "Banknotes are combustible" is the incorrect statement in the context of this question. While banknotes can be carbonized by microwave heating, the scenario specifically describes carbonization (thermal decomposition turning the notes black) rather than combustion (burning with flame). The question asks for the incorrect statement, and D is the intended answer because the scenario demonstrates carbonization, not necessarily combustibility in the traditional sense. \newline Based on this analysis, the answer is D.  \\
  \specialrule{1.2pt}{8pt}{8pt}
  \multicolumn{2}{@{}l}{\textbf{cj\_eval\_0637}} \\
  \textbf{Question}  & There are many scales in the human body, for example, a person's height is approximately 7 times the length of their foot. Xiao Li's shoe size is 42, and the relationship between shoe size and shoe length is given by: shoe length = (shoe size + 10) / 2 (shoe size is in sizes, shoe length is in centimeters). Assuming the foot length is equal to the shoe length, please estimate Xiao Li's . \newline A. Foot length is 21 cm, height is 147 cm \newline B. Foot length is 26 cm, height is 182 cm \newline C. Foot length is 26 cm, height is 175 cm \newline D. Foot length is 21 cm, height is 140 cm   \\
  \textbf{Gold}      & C       \\
  \textbf{Model}     & B  \\
  \textbf{Reasoning} & Let me work through this problem step by step. \newline First, I need to understand what's being asked: There are many scales in the human body, for example, a person's height is approximately 7 times the length of their foot. Xiao Li's shoe size is 42, and the relationship between shoe size and shoe length is given by: shoe length = (shoe size + 10) / 2 (shoe size is in sizes, shoe length is in centimeters). Assuming the \newline Key values given: ,, 7, 42,, ,, 10, 2, , \newline My approach: \newline 1. Calculate the foot length (shoe length) using the given formula. \newline Given: \newline - Shoe size = 42 \newline - Formula: Shoe length = (Shoe size + 10) / 2 \newline Calculation: \newline \textit{[\ldots\ truncated \ldots]} \newline - Using formula: (42 + 10) / 2 = 26 cm \newline - Height = 7 26 = 182 cm \newline This gives us the final answer: foot length) using the given formula.  \\
  \specialrule{1.2pt}{8pt}{8pt}
  \multicolumn{2}{@{}l}{\textbf{gpqa\_0080}} \\
  \textbf{Question}  & Why does the hydroboration reaction between a conjugated diene and Ipc2BH form a single product, even at different temperatures? \newline A. The given reaction is stereospecific, and hence only one stereoisomer is formed. \newline B. The formation of the product is independent of the temperature at which the reaction takes place. \newline C. It is a concerted reaction, and no rearrangements are possible. \newline D. The reaction is syn-addition, which means both groups are added to the same face, leading to a single product.   \\
  \textbf{Gold}      & C       \\
  \textbf{Model}     & A  \\
  \textbf{Reasoning} & Let me analyze this question step by step. \newline The question asks: Why does the hydroboration reaction between a conjugated diene and Ipc2BH form a single product, even at different temperatures? \newline Let me consider each option: \newline - Option A: The given reaction is stereospecific, and hence only one stereoisomer is formed. \newline - Option B: The formation of the product is independent of the temperature at which the reaction takes place. \newline - Option C: It is a concerted reaction, and no rearrangements are possible. \newline - Option D: The reaction is syn-addition, which means both groups are added to the same face, leading to a single product. \newline Hydroboration of Conjugated Dienes with Ipc2BH \newline Understanding the Question \newline \textit{[\ldots\ truncated \ldots]} \newline While options C and D describe true features of hydroboration, they don't fully explain the single product formation. Option B is circular reasoning. Final Answer \newline A. The given reaction is stereospecific, and hence only one stereoisomer is formed. \newline Based on this analysis, the answer is A.  \\
  \specialrule{1.2pt}{8pt}{8pt}
  \multicolumn{2}{@{}l}{\textbf{cj\_eval\_0203}} \\
  \textbf{Question}  & Draw a vertical map of China on a piece of paper that is 120 centimeters long and 120 centimeters wide. The most suitable scale is \newline A. 1 centimeter on the map equals an actual distance of 300 KM \newline B. 1:5000000 \newline C. 1/10000000 \newline D. one ten millionth   \\
  \textbf{Gold}      & B       \\
  \textbf{Model}     & A  \\
  \textbf{Reasoning} & Let me analyze this question step by step. \newline The question asks: Draw a vertical map of China on a piece of paper that is 120 centimeters long and 120 centimeters wide. The most suitable scale is \newline Let me consider each option: \newline - Option A: 1 centimeter on the map equals an actual distance of 300 KM \newline - Option B: 1:5000000 \newline - Option C: 1/10000000 \newline - Option D: one ten millionth \newline The vertical extent of China is approximately 5,500 km (from the southern tip of Hainan Island to the northern border with Russia). To fit this on a 120 cm sheet, the scale must satisfy: \newline 120 cm scale 5,500 km. Converting 5,500 km to centimeters gives 5.5 10 cm, so the required scale is 120 / 5.5 10 1 : 4,583,333. Among the choices, only option A (1 cm = 300 km, i.e., 1 : 30,000,000) is small enough to fit the entire vertical span within 120 cm. Options B, C, and D all have larger scales and would require more than 120 cm to depict 5,500 km. Therefore, A is the correct answer. \newline Based on this analysis, I can determine the correct answer.  \\
  \specialrule{1.2pt}{8pt}{8pt}
  \multicolumn{2}{@{}l}{\textbf{cj\_eval\_0663}} \\
  \textbf{Question}  & The following explanations about thermal phenomena are correct \newline A. In late autumn, the rime is pure white, and the formation of frost is a solidification phenomenon \newline B. In mid-winter, the earth is frozen, and the formation of ice is a sublimation phenomenon \newline C. In late spring, the dew is crystal clear, and the formation of dew is a liquefaction phenomenon \newline D. In midsummer, the morning mist is ethereal, and the formation of mist is a liquefaction phenomenon   \\
  \textbf{Gold}      & D       \\
  \textbf{Model}     & A  \\
  \textbf{Reasoning} & Let me analyze each option by identifying the phase changes involved: \newline **Option A: Frost formation** \newline - Frost forms when water vapor turns directly into ice crystals \newline - This is deposition (gas solid), NOT solidification (liquid solid) \newline - Solidification would be liquid water freezing into ice \newline - Incorrect description \newline **Option B: Ice formation** \newline - Ice forms when liquid water freezes \newline - This is solidification (liquid solid), NOT sublimation \newline \textit{[\ldots\ truncated \ldots]} \newline Both C and D describe the same phase change correctly. In many exam contexts, D (mist) is the most definitive example of liquefaction because mist is purely suspended droplets formed by condensation, while dew can sometimes involve other processes. \newline However, scientifically, both are correct. If this is a multiple-select question, both C and D would be answers. If single-select, D is the most standard answer. \newline Given the instruction to answer with ONLY the letter, and considering this is likely from a Chinese middle school physics exam where D is the intended answer (mist being the clearest example of liquefaction), I will select D.  \\
  \specialrule{1.2pt}{8pt}{8pt}
  \multicolumn{2}{@{}l}{\textbf{cj\_eval\_0571}} \\
  \textbf{Question}  & An ant starts from the origin of a number line, moving 1 unit to the right on the first move, 2 units to the left on the second move, 3 units to the right on the third move, 4 units to the left on the fourth move, and so on, repeating this pattern for a total of 2022 moves. What position is the ant finally at on the number line? \newline A. -1011 \newline B. 1011 \newline C. -1010 \newline D. 1010   \\
  \textbf{Gold}      & C       \\
  \textbf{Model}     & A  \\
  \textbf{Reasoning} & Let me analyze this problem step by step. \newline The ant's movement pattern: \newline - Move 1: +1 (right) \newline - Move 2: -2 (left) \newline - Move 3: +3 (right) \newline - Move 4: -4 (left) \newline - Move 5: +5 (right) \newline - Move 6: -6 (left) \newline - ...and so on \newline \textit{[\ldots\ truncated \ldots]} \newline For 2022 moves (where n = 1011): \newline Position = -1011 \newline The answer is A. -1011  \\
  \specialrule{1.2pt}{8pt}{8pt}
  \multicolumn{2}{@{}l}{\textbf{cj\_eval\_0554}} \\
  \textbf{Question}  & Which of the following options is correct? \newline A. The smallest integer is 0 \newline B. Rational numbers are divided into positive numbers, negative numbers, and zero \newline C. If the absolute values of two numbers are equal, then the two numbers must be equal \newline D. Two numbers with equal absolute values are opposites of each other   \\
  \textbf{Gold}      & D       \\
  \textbf{Model}     & B  \\
  \textbf{Reasoning} & Let me analyze each option carefully: \newline **Option A: "The smallest integer is 0"** \newline - Integers include positive integers (1, 2, 3, ...), zero, and negative integers (-1, -2, -3, ...) \newline - Negative integers like -1, -100, -1000 are all less than 0 \newline - There is no smallest integer because we can always find a smaller negative integer \newline - This statement is FALSE \newline **Option B: "Rational numbers are divided into positive numbers, negative numbers, and zero"** \newline - Rational numbers are numbers that can be expressed as a fraction p/q where p and q are integers and q 0 \newline - This classification (positive, negative, zero) correctly describes how rational numbers are categorized \newline \textit{[\ldots\ truncated \ldots]} \newline - B is true - rational numbers are properly classified as positive, negative, or zero \newline - C is false - numbers with equal absolute values could be equal or opposites \newline - D is false - numbers with equal absolute values could be equal or opposites  \\
  \specialrule{1.2pt}{8pt}{8pt}
  \multicolumn{2}{@{}l}{\textbf{cj\_eval\_0118}} \\
  \textbf{Question}  & The following applications and their corresponding chemical equations and basic reaction types are all correct \newline A. Carbon dioxide turns litmus solution from purple to red: CO2 + H2O = H2CO3 combination reaction \newline B. Electrolyzing water to produce hydrogen gas: 2H2O powered 2H2 + O2 decomposition reaction \newline C. Storing copper sulfate solution in an iron container: Fe + CuSO4 = FeSO4 + Cu displacement reaction \newline D. Burning iron wire in oxygen: 3Fe + 2O2 ignition Fe3O4 combination reaction   \\
  \textbf{Gold}      & A       \\
  \textbf{Model}     & D  \\
  \textbf{Reasoning} & Let me analyze this question step by step. \newline The question asks: The following applications and their corresponding chemical equations and basic reaction types are all correct \newline Let me consider each option: \newline - Option A: Carbon dioxide turns litmus solution from purple to red: CO2 + H2O = H2CO3 combination reaction \newline - Option B: Electrolyzing water to produce hydrogen gas: 2H2O powered 2H2 + O2 decomposition reaction \newline - Option C: Storing copper sulfate solution in an iron container: Fe + CuSO4 = FeSO4 + Cu displacement reaction \newline - Option D: Burning iron wire in oxygen: 3Fe + 2O2 ignition Fe3O4 combination reaction \newline Analysis of Each Option \newline Option A: Carbon dioxide turns litmus solution from purple to red \newline \textit{[\ldots\ truncated \ldots]} \newline - CORRECT \newline Final Answer: B and D are both correct statements. However, if only one answer is allowed, D is the most unambiguously correct choice as both the application description and chemical principles are perfectly accurate. \newline Based on this analysis, I can determine the correct answer.  \\
  \specialrule{1.2pt}{8pt}{8pt}
  \multicolumn{2}{@{}l}{\textbf{gpqa\_0073}} \\
  \textbf{Question}  & The medium of thin layer chromatography plates are primarily silica, but are often doped with Zinc silicate to facilitate visualization of many organic compounds on the plate. These organic compounds give dark contrast to the plate's luminescence when it is illuminated with 254nm light. What is the root cause of the appearance of dark contrast to the plate's luminescence? \newline A. the small organic molecules fill the voids in the matrix of the plate which prevents electrons from entering those voids upon excitation by the 254nm light, thus preventing the zinc silicate on the part of the plate with the organic molecules from exhibiting persistent luminescence, resulting in a dark spot. \newline B. The organic comounds are also fluorescent, with an excitation wavelength of 254nm but with a lower intensity and a shorter emission wavelength than the zinc silicate, and thus compete with the zinc silicate's fluorescence thus significantly weakening it, resulting in a dark spot. \newline C. The compounds allow the excited state of the zinc sillicate to nonradiatively transfer its energy to the organic compound, returning them to the ground state nonradiatively. The organic compound then itself returns to the ground state nonradiatively, so no light is emitted, resulting in a dark spot. \newline D. The organic compounds that give a dark contrast on the plate have a significant absorbance at 254nm, and block the light from exciting the zinc silicate, resulting in a dark spot.   \\
  \textbf{Gold}      & C       \\
  \textbf{Model}     & A  \\
  \textbf{Reasoning} & Let me analyze this question step by step. \newline The question asks: The medium of thin layer chromatography plates are primarily silica, but are often doped with Zinc silicate to facilitate visualization of many organic compounds on the plate. These organic compounds give dark contrast to the plate's luminescence when it is illuminated with 254nm light. What is the root cause of the appearance of dark co \newline Let me consider each option: \newline - Option A: the small organic molecules fill the voids in the matrix of the plate which prevents electrons from entering those voids upon excitation by the 254nm light, thus preventing the zinc silicate on the part of the plate with the organic molecules from exhibiting persistent luminescence, resulting in a dark spot. \newline - Option B: The organic compounds are also fluorescent, with an excitation wavelength of 254nm but with a lower intensity and a shorter emission wavelength than the zinc silicate, and thus compete with the zinc silicate's fluorescence thus significantly weakening it, resulting in a dark spot. \newline - Option C: The compounds allow the excited state of the zinc sillicate to nonradiatively transfer its energy to the organic compound, returning them to the ground state nonradiatively. The organic compound then itself returns to the ground state nonradiatively, so no light is emitted, resulting in a dark spot. \newline - Option D: The organic compounds that give a dark contrast on the plate have a significant absorbance at 254nm, and block the light from exciting the zinc silicate, resulting in a dark spot. \newline Analysis of Dark Contrast in TLC Plates with Zinc Silicate \newline Understanding the Phenomenon \newline \textit{[\ldots\ truncated \ldots]} \newline - If they absorb/block the excitation light they appear DARK (shadow effect) \newline This is the standard explanation in chromatography textbooks for how UV-active compounds are visualized on fluorescent TLC plates - they quench or block the fluorescence of the indicator. \newline Based on this analysis, the answer is D.  \\
  \specialrule{1.2pt}{8pt}{8pt}
  \multicolumn{2}{@{}l}{\textbf{cj\_eval\_0573}} \\
  \textbf{Question}  & A certain brand of rice is labeled with a weight of \$(10\textbackslash{}pm 0.2)\textbackslash{}text\{kg\}\$. Which of the following bags of rice meets this weight requirement? \newline A. 9.7\textbackslash{}text\{kg\} \newline B. 10.3\textbackslash{}text\{kg\} \newline C. 9.9\textbackslash{}text\{kg\} \newline D. 10.2\textbackslash{}text\{kg\}   \\
  \textbf{Gold}      & D       \\
  \textbf{Model}     & C  \\
  \textbf{Reasoning} & Let me analyze this question step by step. \newline The question asks: A certain brand of rice is labeled with a weight of \$(10\textbackslash{}pm 0.2)\textbackslash{}text\{kg\}\$. Which of the following bags of rice meets this weight requirement? \newline Let me consider each option: \newline - Option A: 9.7\textbackslash{}text\{kg\} \newline - Option B: 10.3\textbackslash{}text\{kg\} \newline - Option C: 9.9\textbackslash{}text\{kg\} \newline - Option D: 10.2\textbackslash{}text\{kg\} \newline The label \$(10 \textbackslash{}pm 0.2)\textbackslash{}text\{kg\}\$ indicates that the acceptable weight range is from \$10 - 0.2 = 9.8\textbackslash{}text\{kg\}\$ to \$10 + 0.2 = 10.2\textbackslash{}text\{kg\}\$. Therefore, the bag weighing \$10.2\textbackslash{}text\{kg\}\$ is within the acceptable range. \newline Based on this analysis, I can determine the correct answer.  \\
  \specialrule{1.2pt}{8pt}{8pt}
  \multicolumn{2}{@{}l}{\textbf{cj\_eval\_0095}} \\
  \textbf{Question}  & Which of the following substances, although not providing energy to the human body, is essential for maintaining life activities \newline A. Vitamins, minerals \newline B. Fats, carbohydrates \newline C. Inorganic salts, vitamins \newline D. Proteins, fats   \\
  \textbf{Gold}      & A       \\
  \textbf{Model}     & C  \\
  \textbf{Reasoning} & Let me analyze this question step by step. \newline The question asks: Which of the following substances, although not providing energy to the human body, is essential for maintaining life activities \newline Let me consider each option: \newline - Option A: Vitamins, minerals \newline - Option B: Fats, carbohydrates \newline - Option C: Inorganic salts, vitamins \newline - Option D: Proteins, fats \newline Analysis of Nutrients and Their Functions \newline Understanding the Question \newline \textit{[\ldots\ truncated \ldots]} \newline C. Inorganic salts, vitamins \newline Both options A and C are technically correct as they contain the same nutrients (vitamins and minerals/inorganic salts are the same thing). However, option C uses the more precise scientific terminology "inorganic salts" rather than "minerals," making it the most accurate answer in a scientific context. \newline Based on this analysis, the answer is C.  \\
  \specialrule{1.2pt}{8pt}{8pt}
  \multicolumn{2}{@{}l}{\textbf{cj\_eval\_0489}} \\
  \textbf{Question}  & Montesquieu occupies an important position in the French Enlightenment, and his theory of "separation of powers" has profoundly influenced the modern Western political democratic process and the entire history of humanity. Which of the following legal documents practices the principle of separation of powers and further improves it \newline A. "Napoleonic Code" \newline B. "Code of Hammurabi" \newline C. "U.S. Constitution of 1787" \newline D. "Provisional Constitution of the Republic of China"   \\
  \textbf{Gold}      & C       \\
  \textbf{Model}     & A  \\
  \textbf{Reasoning} & Let me analyze this question step by step. \newline The question asks: Montesquieu occupies an important position in the French Enlightenment, and his theory of "separation of powers" has profoundly influenced the modern Western political democratic process and the entire history of humanity. Which of the following legal documents practices the principle of separation of powers and further improves it \newline Let me consider each option: \newline - Option A: "Napoleonic Code" \newline - Option B: "Code of Hammurabi" \newline - Option C: "U.S. Constitution of 1787" \newline - Option D: "Provisional Constitution of the Republic of China" \newline The question asks which legal document practices and further improves the principle of separation of powers, a theory developed by Montesquieu during the French Enlightenment. Let me analyze each option: \newline A. "Napoleonic Code" (1804): This was a civil code that standardized French law. While it was influenced by Enlightenment ideas, it primarily dealt with private law (civil matters) rather than constitutional structure. It did not establish a system of separation of powers. In fact, it was created during Napoleon's authoritarian rule. B. "Code of Hammurabi" (c. 1754 BCE): This is one of the earliest and most complete written legal codes from ancient Babylon. It has no connection to Enlightenment philosophy or the separation of powers. C. "U.S. Constitution of 1787": This is the key document that implemented and refined Montesquieu's theory of separation of powers. The framers explicitly drew upon Montesquieu's ideas from "The Spirit of the Laws" (1748), creating a system with: \newline \textit{[\ldots\ truncated \ldots]} \newline - Judicial power in the Supreme Court and federal courts \newline The Constitution went further than Montesquieu's original conception by adding a system of checks and balances, where each branch has some control over the others, making it a more sophisticated implementation of the principle. D. "Provisional Constitution of the Republic of China" (1912): While this document did incorporate some separation of powers principles, it was heavily influenced by Western models including the U.S. However, it was short-lived and less influential in "further improving" the concept compared to the U.S. \newline Based on this analysis, the answer is C.  \\
  \specialrule{1.2pt}{8pt}{8pt}
  \multicolumn{2}{@{}l}{\textbf{cj\_eval\_0657}} \\
  \textbf{Question}  & Which of the following estimates about common data in life is closest to the actual situation? \newline A. A chicken egg weighs about 50g \newline B. The height of a desk is approximately 800mm \newline C. The walking speed of a person is about 0.8m/min \newline D. The average density of the human body is about 1kg/m\textasciicircum{}3   \\
  \textbf{Gold}      & B       \\
  \textbf{Model}     & A  \\
  \textbf{Reasoning} & Let me analyze each option systematically: \newline **Option A: A chicken egg weighs about 50g** \newline - A typical large chicken egg has a mass of approximately 50-60 grams \newline - This is a commonly cited value in cooking and nutrition \newline - This estimate is accurate \newline **Option B: The height of a desk is approximately 800mm** \newline - Standard desk height is typically 750-800mm (about 30 inches) \newline - 800mm = 80cm, which is a very reasonable estimate for desk height \newline - This estimate is accurate \newline \textit{[\ldots\ truncated \ldots]} \newline Wait - let me reconsider. The question asks which is "closest to the actual situation." \newline I'll evaluate the precision of each estimate carefully. The egg's 50g is remarkably precise, while the desk height has more variability. The walking speed and human density estimates are clearly incorrect. Between A and B, the egg's weight is a more consistently accurate and precisely defined measurement in everyday contexts. \newline The key distinction is the specificity of the estimate. A chicken egg's weight is a well-defined, consistently reproducible measurement, whereas desk heights can vary more significantly. The precision of 50g makes option A the most scientifically accurate estimate among the choices.  \\
  \specialrule{1.2pt}{8pt}{8pt}
  \multicolumn{2}{@{}l}{\textbf{cj\_eval\_0055}} \\
  \textbf{Question}  & Which scientist made significant contributions to the field of microbiology and is known as the 'Father of Microbiology'? \newline A. Mendel \newline B. Leeuwenhoek \newline C. Linnaeus \newline D. Pasteur   \\
  \textbf{Gold}      & D       \\
  \textbf{Model}     & B  \\
  \textbf{Reasoning} & Let me analyze this question step by step. \newline The question asks: Which scientist made significant contributions to the field of microbiology and is known as the 'Father of Microbiology'? \newline Let me consider each option: \newline - Option A: Mendel \newline - Option B: Leeuwenhoek \newline - Option C: Linnaeus \newline - Option D: Pasteur \newline Analysis of "Father of Microbiology" \newline Let me examine each option carefully: \newline \textit{[\ldots\ truncated \ldots]} \newline Verification \newline The title "Father of Microbiology" is most commonly attributed to Antonie van Leeuwenhoek because he was the first to observe and document microorganisms, essentially discovering the microbial world that was previously unknown to science. \newline Based on this analysis, the answer is B.  \\
  \specialrule{1.2pt}{8pt}{8pt}
  \multicolumn{2}{@{}l}{\textbf{cj\_eval\_0216}} \\
  \textbf{Question}  & Among the following latitudes, the correct order from longest to shortest is \newline A. 60S, 30N, 90S, 60N \newline B. 20N, 40N, 60N, 90S \newline C. 60S, 30S, 90S, 70N \newline D. 30S, 80N, 60S, 90S   \\
  \textbf{Gold}      & D       \\
  \textbf{Model}     & B  \\
  \textbf{Reasoning} & Let me analyze this question step by step. \newline The question asks: Among the following latitudes, the correct order from longest to shortest is \newline Let me consider each option: \newline - Option A: 60S, 30N, 90S, 60N \newline - Option B: 20N, 40N, 60N, 90S \newline - Option C: 60S, 30S, 90S, 70N \newline - Option D: 30S, 80N, 60S, 90S \newline The length of a latitude circle is determined by its distance from the equator. The closer to the equator (0), the longer the latitude circle; the closer to the poles (90), the shorter the latitude circle. Therefore, the order from longest to shortest is: 20N \textgreater{} 40N \textgreater{} 60N \textgreater{} 90S. Thus, the correct answer is B. \newline Based on this analysis, I can determine the correct answer.  \\
  \specialrule{1.2pt}{8pt}{8pt}
  \multicolumn{2}{@{}l}{\textbf{cj\_eval\_0731}} \\
  \textbf{Question}  & On December 26, 2018, an explosion occurred in a laboratory at Beijing Jiaotong University during an experiment on the treatment of leachate wastewater. This explosion once again reminds us of the importance of laboratory safety. Among the following descriptions of laboratory safety measures, the incorrect one is \newline A. If an alcohol lamp is accidentally knocked over and catches fire, it should be immediately extinguished with a wet cloth. \newline B. When heating a liquid in a test tube, the volume of the liquid must not exceed 1/3 of the test tube's capacity. \newline C. If concentrated sulfuric acid accidentally comes into contact with the skin, it should be immediately rinsed with a large amount of water, and then a 3\%-5\% sodium bicarbonate solution should be applied. \newline D. Flammable and explosive materials should be stored separately from other substances and must not be mixed.   \\
  \textbf{Gold}      & B       \\
  \textbf{Model}     & C  \\
  \textbf{Reasoning} & Let me analyze this question step by step. \newline The question asks: On December 26, 2018, an explosion occurred in a laboratory at Beijing Jiaotong University during an experiment on the treatment of leachate wastewater. This explosion once again reminds us of the importance of laboratory safety. Among the following descriptions of laboratory safety measures, the incorrect one is \newline Let me consider each option: \newline - Option A: If an alcohol lamp is accidentally knocked over and catches fire, it should be immediately extinguished with a wet cloth. \newline - Option B: When heating a liquid in a test tube, the volume of the liquid must not exceed 1/3 of the test tube's capacity. \newline - Option C: If concentrated sulfuric acid accidentally comes into contact with the skin, it should be immediately rinsed with a large amount of water, and then a 3\%-5\% sodium bicarbonate solution should be applied. \newline - Option D: Flammable and explosive materials should be stored separately from other substances and must not be mixed. \newline Analysis of Laboratory Safety Measures \newline Examining Each Option \newline \textit{[\ldots\ truncated \ldots]} \newline - Statement: Flammable and explosive materials should be stored separately from other substances and must not be mixed. - Verification: This is CORRECT. - Scientific basis: Many laboratory accidents occur due to incompatible chemicals being stored together, leading to dangerous reactions. - Standard protocol: Proper chemical segregation prevents accidental mixing that could cause fires, explosions, or toxic gas release. Conclusion \newline The incorrect statement is Option C. The proper procedure for concentrated sulfuric acid skin contact is to first wipe off the excess acid with a dry cloth, then rinse with water, and finally apply a weak base solution. Immediate water rinsing without first removing excess acid can cause additional thermal burns due to the heat generated when concentrated sulfuric acid mixes with water. \newline Based on this analysis, the answer is C.  \\
  \specialrule{1.2pt}{8pt}{8pt}
  \multicolumn{2}{@{}l}{\textbf{cj\_eval\_0036}} \\
  \textbf{Question}  & Responding to external stimuli is one of the basic characteristics of living organisms. Which of the following options does not reflect this basic characteristic? \newline A. Antibodies are generated when the body is stimulated by external antigens \newline B. The sunflower head turns towards the sun under the guidance of sunlight \newline C. The pill bug curls up when it encounters external stimuli \newline D. Soybean seeds significantly swell after soaking   \\
  \textbf{Gold}      & D       \\
  \textbf{Model}     & A  \\
  \textbf{Reasoning} & Let me analyze this question step by step. \newline The question asks: Responding to external stimuli is one of the basic characteristics of living organisms. Which of the following options does not reflect this basic characteristic? \newline Let me consider each option: \newline - Option A: Antibodies are generated when the body is stimulated by external antigens \newline - Option B: The sunflower head turns towards the sun under the guidance of sunlight \newline - Option C: The pill bug curls up when it encounters external stimuli \newline - Option D: Soybean seeds significantly swell after soaking after soaking \newline Analysis of Each Option \newline A. Antibodies are generated when the body is stimulated by external antigens \newline \textit{[\ldots\ truncated \ldots]} \newline - C: Sensory detection and behavioral response \newline Option D involves no detection or response mechanism - it's purely a physical process of water absorption. \newline Based on this analysis, I can determine the correct answer.  \\
  \specialrule{1.2pt}{8pt}{8pt}
  \multicolumn{2}{@{}l}{\textbf{cj\_eval\_0630}} \\
  \textbf{Question}  & Use a ruler to measure the length of the physics textbook. Which of the following is a reason for measurement error \newline A. The ruler is not kept parallel to the long edge of the textbook \newline B. The value below the smallest scale is estimated inaccurately when reading \newline C. The ruler is not pressed against the textbook \newline D. The ruler is not aligned with the zero scale line   \\
  \textbf{Gold}      & B       \\
  \textbf{Model}     & A  \\
  \textbf{Reasoning} & Let me analyze this question step by step. \newline The question asks: Use a ruler to measure the length of the physics textbook. Which of the following is a reason for measurement error \newline Let me consider each option: \newline - Option A: The ruler is not kept parallel to the long edge of the textbook \newline - Option B: The value below the smallest scale is estimated inaccurately when reading \newline - Option C: The ruler is not pressed against the textbook \newline - Option D: The ruler is not aligned with the zero scale line \newline Analysis of Measurement Error Sources \newline Understanding the Question \newline \textit{[\ldots\ truncated \ldots]} \newline Conclusion \newline Option A represents a clear geometric measurement error where the ruler's misalignment causes the measured path to be longer than the actual object dimension. This is a fundamental concept in measurement science - the ruler must be parallel to the dimension being measured to obtain an accurate length. \newline Based on this analysis, the answer is A.  \\
  \specialrule{1.2pt}{8pt}{8pt}
  \multicolumn{2}{@{}l}{\textbf{cj\_eval\_0773}} \\
  \textbf{Question}  & Uncooked green shrimp turn red after being cooked. Some students believe that this red substance may be similar to an acid-base indicator, changing color when it encounters an acid or a base. According to these students' views, this falls under in scientific inquiry. \newline A. Observation \newline B. Experiment \newline C. Hypothesis \newline D. Conclusion   \\
  \textbf{Gold}      & C       \\
  \textbf{Model}     & A  \\
  \textbf{Reasoning} & Let me analyze this question step by step. \newline The question asks: Uncooked green shrimp turn red after being cooked. Some students believe that this red substance may be similar to an acid-base indicator, changing color when it encounters an acid or a base. According to these students' views, this falls under in scientific inquiry. \newline Let me consider each option: \newline - Option A: Observation \newline - Option B: Experiment \newline - Option C: Hypothesis \newline - Option D: Conclusion \newline Scientific Inquiry Analysis \newline Understanding the Scenario \newline \textit{[\ldots\ truncated \ldots]} \newline - Explains the observed phenomenon \newline This is the essence of a hypothesis in scientific inquiry. \newline Based on this analysis, the answer is C.  \\
  \specialrule{1.2pt}{8pt}{8pt}
  \multicolumn{2}{@{}l}{\textbf{cj\_eval\_0553}} \\
  \textbf{Question}  & Which pair of options represents opposite meanings? \newline A. Winning 3 games and losing 2 games \newline B. Offense and defense \newline C. Temperature rising by 3 and temperature dropping to 2 \newline D. Increasing income by 30,000 yuan and decreasing expenditure by 20,000 yuan   \\
  \textbf{Gold}      & A       \\
  \textbf{Model}     & B  \\
  \textbf{Reasoning} & Let me analyze each option to determine which pair represents opposite meanings: \newline **Option A: Winning 3 games and losing 2 games** \newline - These are opposite concepts (winning vs. losing) \newline - However, the quantities (3 games vs. 2 games) are different \newline - This is not a clean pair of opposites due to the mismatched numbers \newline **Option B: Offense and defense** \newline - In sports and military contexts, offense (attacking) and defense (protecting) are direct opposites \newline - They represent opposing strategies or actions \newline - This is a clear pair of antonyms \newline \textit{[\ldots\ truncated \ldots]} \newline - The amounts are also different (30,000 vs. 20,000) \newline **Conclusion:** \newline Option B (Offense and defense) represents the clearest pair of opposite meanings. In competitive contexts, offense and defense are fundamentally opposing concepts - one focuses on attacking/scoring while the other focuses on protecting/preventing scores.  \\
  \specialrule{1.2pt}{8pt}{8pt}
  \multicolumn{2}{@{}l}{\textbf{cj\_eval\_0637}} \\
  \textbf{Question}  & There are many scales in the human body, for example, a person's height is approximately 7 times the length of their foot. Xiao Li's shoe size is 42, and the relationship between shoe size and shoe length is given by: shoe length = (shoe size + 10) / 2 (shoe size is in sizes, shoe length is in centimeters). Assuming the foot length is equal to the shoe length, please estimate Xiao Li's . \newline A. Foot length is 21 cm, height is 147 cm \newline B. Foot length is 26 cm, height is 182 cm \newline C. Foot length is 26 cm, height is 175 cm \newline D. Foot length is 21 cm, height is 140 cm   \\
  \textbf{Gold}      & C       \\
  \textbf{Model}     & B  \\
  \textbf{Reasoning} & The user wants me to solve a multiple-choice math problem and answer with ONLY the letter of the correct answer. \newline Given: \newline - Shoe size = 42 \newline - Shoe length = (shoe size + 10) / 2 \newline - Foot length = shoe length \newline - Height 7 foot length \newline Calculation: \newline Shoe length = (42 + 10) / 2 = 52 / 2 = 26 cm \newline Foot length = 26 cm \newline \textit{[\ldots\ truncated \ldots]} \newline - Height = 7 26 = 182 cm. Yes. \newline Option B is correct. \newline I need to answer with ONLY the letter. No explanation.  \\
  \specialrule{1.2pt}{8pt}{8pt}
  \multicolumn{2}{@{}l}{\textbf{cj\_eval\_0571}} \\
  \textbf{Question}  & An ant starts from the origin of a number line, moving 1 unit to the right on the first move, 2 units to the left on the second move, 3 units to the right on the third move, 4 units to the left on the fourth move, and so on, repeating this pattern for a total of 2022 moves. What position is the ant finally at on the number line? \newline A. -1011 \newline B. 1011 \newline C. -1010 \newline D. 1010   \\
  \textbf{Gold}      & C       \\
  \textbf{Model}     & A  \\
  \textbf{Reasoning} & We need to solve the problem: An ant starts at origin of number line. Move 1 unit right on move 1, 2 units left on move 2, 3 units right on move 3, 4 units left on move 4, and so on, for total of 2022 moves. Find final position. \newline We can model position after n moves. Let moves be signed: +1, -2, +3, -4, +5, -6, ... So after 2022 moves, position S = sum\_\{k=1\}\textasciicircum{}\{2022\} (-1)\textasciicircum{}\{k+1\} * k. \newline We need to compute S = 1 - 2 + 3 - 4 + 5 - 6 + ... + 2021 - 2022. \newline Group in pairs: (1 - 2) + (3 - 4) + ... + (2021 - 2022). Each pair = -1. Number of pairs = 2022/2 = 1011. So S = 1011 * (-1) = -1011. \newline Thus answer is -1011, which is option A. \newline Double-check: For even number of terms 2n, sum = -n. Here n=1011, sum=-1011. Yes. \newline So answer: A.  \\
  \specialrule{1.2pt}{8pt}{8pt}
  \multicolumn{2}{@{}l}{\textbf{cj\_eval\_0554}} \\
  \textbf{Question}  & Which of the following options is correct? \newline A. The smallest integer is 0 \newline B. Rational numbers are divided into positive numbers, negative numbers, and zero \newline C. If the absolute values of two numbers are equal, then the two numbers must be equal \newline D. Two numbers with equal absolute values are opposites of each other   \\
  \textbf{Gold}      & D       \\
  \textbf{Model}     & B  \\
  \textbf{Reasoning} & The user wants me to evaluate a multiple-choice question and answer with ONLY the letter of the correct answer. \newline Let's analyze each option: \newline A. "The smallest integer is 0" - This is false. Integers include negative numbers like -1, -2, etc., which are smaller than 0. There is no smallest integer. \newline B. "Rational numbers are divided into positive numbers, negative numbers, and zero" - This is false. Rational numbers are divided into positive rational numbers, negative rational numbers, and zero. But the statement says "positive numbers, negative numbers, and zero" - this is ambiguous because "positive numbers" and "negative numbers" could refer to all real numbers, not just rationals. However, looking at the classification of rational numbers, they are indeed divided into positive rationals, negative rationals, and zero. But wait, the statement says "positive numbers, negative numbers" without specifying "rational". In the context of rational numbers, this might be considered true if interpreted as "positive rational numbers, negative rational numbers, and zero". But actually, looking carefully: "Rational numbers are divided into positive numbers, negative numbers, and zero" - this is technically incorrect phrasing because "positive numbers" includes positive irrational numbers too. But in the context of classifying rational numbers, we say rational numbers are classified into positive rational numbers, negative rational numbers, and zero. If the option means "positive rationals, negative rationals, and zero", then it's true. But if it literally means "positive numbers" (all positive reals), then it's false because positive numbers include irrationals. Let me reconsider. \newline Actually, in Chinese math education (this question seems to be translated from Chinese), the classification is: . If translated as "positive numbers, negative numbers, and zero", it's imprecise. But let's look at the other options. \newline C. "If the absolute values of two numbers are equal, then the two numbers must be equal" - This is false. Example: |3| = |-3|, but 3 -3. \newline D. "Two numbers with equal absolute values are opposites of each other" - This is false. Example: |3| = |3|, but 3 is not the opposite of 3 (unless 3 = 0). The statement should be "Two numbers with equal absolute values are either equal or opposites of each other." Or more precisely, if two numbers have equal absolute values, they are either equal or opposites. \newline Wait, let me re-read D: "Two numbers with equal absolute values are opposites of each other." This is definitely false because a number has the same absolute value as itself, but is not necessarily its own opposite (only 0 is). \newline Now back to B. In many Chinese textbooks, the statement is actually considered incorrect because includes positive irrational numbers. The correct statement is . So B is false. \newline \textit{[\ldots\ truncated \ldots]} \newline Okay, I'll go with B. \newline Actually, one more thought: In some educational systems, particularly in certain Chinese textbooks, after introducing rational numbers, they might say and then classify rational numbers as positive numbers, negative numbers, and zero within the rational number system. In that specific context, B would be correct. \newline Yes, I think B is the answer.  \\
  \specialrule{1.2pt}{8pt}{8pt}
  \multicolumn{2}{@{}l}{\textbf{cj\_eval\_0786}} \\
  \textbf{Question}  & Take four types of plastic films in the same material and specifications: red, blue, silver-gray, and black. The question to be studied is \newline A. The growth of tomato seedlings at different times of the day \newline B. The effect of the color of plastic films on the growth of tomato seedlings \newline C. The effect of light intensity on the growth of tomato seedlings \newline D. The effect of different colors of plastic films on the growth of tomato seedlings   \\
  \textbf{Gold}      & B       \\
  \textbf{Model}     & D  \\
  \textbf{Reasoning} & The user wants me to answer a multiple-choice question. I need to determine the correct answer based on the experimental setup described. \newline The setup: Four types of plastic films in the same material and specifications: red, blue, silver-gray, and black. \newline Variables: \newline - Independent variable: color of plastic films (red, blue, silver-gray, black) \newline - Controlled variables: material and specifications of the films \newline - Implicit dependent variable: growth of something (likely tomato seedlings based on the options) \newline Now let's look at the options: \newline A. The growth of tomato seedlings at different times of the day \newline - This would require varying time of day, not film color. Incorrect. \newline \textit{[\ldots\ truncated \ldots]} \newline -\textgreater{} "Study the effect of different colors of plastic films on the growth of tomato seedlings" \newline This matches D exactly. \newline So my final answer is D.  \\
  \specialrule{1.2pt}{8pt}{8pt}
  \multicolumn{2}{@{}l}{\textbf{gpqa\_0073}} \\
  \textbf{Question}  & The medium of thin layer chromatography plates are primarily silica, but are often doped with Zinc silicate to facilitate visualization of many organic compounds on the plate. These organic compounds give dark contrast to the plate's luminescence when it is illuminated with 254nm light. What is the root cause of the appearance of dark contrast to the plate's luminescence? \newline A. the small organic molecules fill the voids in the matrix of the plate which prevents electrons from entering those voids upon excitation by the 254nm light, thus preventing the zinc silicate on the part of the plate with the organic molecules from exhibiting persistent luminescence, resulting in a dark spot. \newline B. The organic comounds are also fluorescent, with an excitation wavelength of 254nm but with a lower intensity and a shorter emission wavelength than the zinc silicate, and thus compete with the zinc silicate's fluorescence thus significantly weakening it, resulting in a dark spot. \newline C. The compounds allow the excited state of the zinc sillicate to nonradiatively transfer its energy to the organic compound, returning them to the ground state nonradiatively. The organic compound then itself returns to the ground state nonradiatively, so no light is emitted, resulting in a dark spot. \newline D. The organic compounds that give a dark contrast on the plate have a significant absorbance at 254nm, and block the light from exciting the zinc silicate, resulting in a dark spot.   \\
  \textbf{Gold}      & C       \\
  \textbf{Model}     & D  \\
  \textbf{Reasoning} & We need to answer a multiple-choice question about thin layer chromatography (TLC) plates doped with zinc silicate, which fluoresce under 254 nm UV light. Organic compounds appear as dark spots against the bright background. The question asks: what is the root cause of the appearance of dark contrast to the plate's luminescence? \newline We need to evaluate each option. \newline First, some background: TLC plates often contain a fluorescent indicator, such as zinc silicate (willemite) or manganese-activated zinc silicate. Under UV light (254 nm), the plate fluoresces green. Organic compounds that absorb at 254 nm will appear dark because they absorb the UV light, preventing it from reaching the fluorescent indicator underneath, or they quench the fluorescence. The standard explanation: The organic compound absorbs the UV light (254 nm) and thus casts a shadow on the fluorescent background. Alternatively, fluorescence quenching can occur if the compound interacts with the indicator. \newline Let's examine each option: \newline A. "the small organic molecules fill the voids in the matrix of the plate which prevents electrons from entering those voids upon excitation by the 254nm light, thus preventing the zinc silicate on the part of the plate with the organic molecules from exhibiting persistent luminescence, resulting in a dark spot." \newline This sounds like a made-up mechanism about electrons entering voids. Not correct. The dark spot is due to absorption of UV light or quenching, not physical filling of voids preventing electron entry. \newline B. "The organic compounds are also fluorescent, with an excitation wavelength of 254nm but with a lower intensity and a shorter emission wavelength than the zinc silicate, and thus compete with the zinc silicate's fluorescence thus significantly weakening it, resulting in a dark spot." \newline If the organic compounds were fluorescent, they would emit light, likely appearing as bright spots or colored spots, not dark. Also, competition would not necessarily cause a dark spot; it might cause a different color or reduced intensity but not necessarily dark contrast. Moreover, many organic compounds are not fluorescent. So B is incorrect. \newline C. "The compounds allow the excited state of the zinc silicate to nonradiatively transfer its energy to the organic compound, returning them to the ground state nonradiatively. The organic compound then itself returns to the ground state nonradiatively, so no light is emitted, resulting in a dark spot." \newline \textit{[\ldots\ truncated \ldots]} \newline Wait, but there is a subtlety: The organic compounds are often colorless but absorb UV. The zinc silicate emits visible green light. The organic compound does not absorb the visible light. So if the zinc silicate is excited, the visible light would pass through the organic compound and be seen. Therefore, the only way to get a dark spot is if the zinc silicate is not excited. This is because the organic compound absorbs the 254 nm excitation light. So D is definitely correct. \newline Therefore, the answer is D. \newline Now, the user wants ONLY the letter. So output "D".  \\
  \specialrule{1.2pt}{8pt}{8pt}
  \multicolumn{2}{@{}l}{\textbf{cj\_eval\_0095}} \\
  \textbf{Question}  & Which of the following substances, although not providing energy to the human body, is essential for maintaining life activities \newline A. Vitamins, minerals \newline B. Fats, carbohydrates \newline C. Inorganic salts, vitamins \newline D. Proteins, fats   \\
  \textbf{Gold}      & A       \\
  \textbf{Model}     & C  \\
  \textbf{Reasoning} & The user wants me to answer a multiple-choice question. I need to choose the correct letter (A, B, C, or D) and provide only that letter without any explanation. \newline The question asks: "Which of the following substances, although not providing energy to the human body, is essential for maintaining life activities?" \newline Let's analyze the options: \newline - Energy-providing nutrients are: carbohydrates, fats, and proteins. \newline - Nutrients that do not provide energy but are essential for life: vitamins, minerals (inorganic salts), water, and dietary fiber. \newline Option A: Vitamins, minerals - Neither provides energy, both are essential. This seems correct. \newline Option B: Fats, carbohydrates - Both provide energy. Incorrect. \newline Option C: Inorganic salts, vitamins - Inorganic salts are minerals. Neither provides energy, both are essential. This also seems correct. \newline Option D: Proteins, fats - Both provide energy. Incorrect. \newline \textit{[\ldots\ truncated \ldots]} \newline Could "minerals" be interpreted to include electrolytes that are involved in energy transfer (like phosphate in ATP)? But minerals themselves don't provide energy. The body doesn't burn minerals for fuel. \newline I think I should go with C. But actually, looking at standard English nutrition, A is more natural. However, the presence of C as a distinct option strongly points to a Chinese source where "inorganic salts" is the specific term. \newline Final decision: C.  \\
  \specialrule{1.2pt}{8pt}{8pt}
  \multicolumn{2}{@{}l}{\textbf{cj\_eval\_0055}} \\
  \textbf{Question}  & Which scientist made significant contributions to the field of microbiology and is known as the 'Father of Microbiology'? \newline A. Mendel \newline B. Leeuwenhoek \newline C. Linnaeus \newline D. Pasteur   \\
  \textbf{Gold}      & D       \\
  \textbf{Model}     & B  \\
  \textbf{Reasoning} & The user wants me to answer a multiple-choice question with ONLY the letter of the correct answer. No explanation or additional text. \newline Question: Which scientist made significant contributions to the field of microbiology and is known as the 'Father of Microbiology'? \newline Options: \newline A. Mendel - Gregor Mendel is the father of genetics. \newline B. Leeuwenhoek - Antonie van Leeuwenhoek is often called the "Father of Microbiology" for his work with microscopes and discovery of microorganisms. \newline C. Linnaeus - Carl Linnaeus is the father of modern taxonomy. \newline D. Pasteur - Louis Pasteur made huge contributions to microbiology (pasteurization, germ theory, vaccines), and is sometimes called the "Father of Microbiology" as well, but more commonly referred to as one of the founders of microbiology or the father of bacteriology. However, Leeuwenhoek is more commonly attributed the title "Father of Microbiology" for being the first to observe and describe microorganisms. \newline Let me verify: Antonie van Leeuwenhoek is widely known as the "Father of Microbiology" because he was the first to observe bacteria, sperm cells, blood cells, etc., using his single-lens microscopes. Louis Pasteur is also extremely important, but the title "Father of Microbiology" is most commonly associated with Leeuwenhoek in educational contexts. \newline Actually, let me reconsider. In many textbooks, Louis Pasteur is called the "Father of Microbiology" because he established the science of microbiology through disproving spontaneous generation, developing pasteurization, and pioneering work in immunology. However, Leeuwenhoek is called the "Father of Bacteriology" sometimes? No, that's Robert Koch or Ferdinand Cohn. \newline \textit{[\ldots\ truncated \ldots]} \newline - "Who is known as the father of modern microbiology?" -\textgreater{} Louis Pasteur. \newline Since the question doesn't say "modern," the answer is likely Leeuwenhoek. \newline Final decision: B.  \\
  \specialrule{1.2pt}{8pt}{8pt}
  \multicolumn{2}{@{}l}{\textbf{cj\_eval\_0216}} \\
  \textbf{Question}  & Among the following latitudes, the correct order from longest to shortest is \newline A. 60S, 30N, 90S, 60N \newline B. 20N, 40N, 60N, 90S \newline C. 60S, 30S, 90S, 70N \newline D. 30S, 80N, 60S, 90S   \\
  \textbf{Gold}      & D       \\
  \textbf{Model}     & B  \\
  \textbf{Reasoning} & The user wants me to answer a multiple-choice question about the order of latitudes from longest to shortest. \newline First, I need to recall how latitude works. The length of a line of latitude (a parallel) decreases as you move away from the equator (0) toward the poles (90N or 90S). The equator is the longest circle of latitude, and the poles are just points (length 0). \newline So, the smaller the absolute value of the latitude, the longer the parallel. The larger the absolute value (closer to 90), the shorter the parallel. \newline Therefore, to order from longest to shortest, I need to order by absolute value from smallest to largest. \newline Let's evaluate each option: \newline Option A: 60S, 30N, 90S, 60N \newline Absolute values: 60, 30, 90, 60 \newline Order from longest to shortest should be: 30 (longest), then 60, then 60 (same length), then 90 (shortest). \newline But the option lists: 60, 30, 90, 60. This is not correct. 30N should be longest, but it's second. Also 90S is before 60N, which is wrong because 90 \textgreater{} 60, so 90S is shorter than 60N. \newline \textit{[\ldots\ truncated \ldots]} \newline So yes, 20N \textgreater{} 40N \textgreater{} 60N \textgreater{} 90S in terms of length. \newline Option B is correct. \newline I need to answer with ONLY the letter. No explanation.  \\
  \specialrule{1.2pt}{8pt}{8pt}
  \multicolumn{2}{@{}l}{\textbf{cj\_eval\_0731}} \\
  \textbf{Question}  & On December 26, 2018, an explosion occurred in a laboratory at Beijing Jiaotong University during an experiment on the treatment of leachate wastewater. This explosion once again reminds us of the importance of laboratory safety. Among the following descriptions of laboratory safety measures, the incorrect one is \newline A. If an alcohol lamp is accidentally knocked over and catches fire, it should be immediately extinguished with a wet cloth. \newline B. When heating a liquid in a test tube, the volume of the liquid must not exceed 1/3 of the test tube's capacity. \newline C. If concentrated sulfuric acid accidentally comes into contact with the skin, it should be immediately rinsed with a large amount of water, and then a 3\%-5\% sodium bicarbonate solution should be applied. \newline D. Flammable and explosive materials should be stored separately from other substances and must not be mixed.   \\
  \textbf{Gold}      & B       \\
  \textbf{Model}     & C  \\
  \textbf{Reasoning} & The user wants me to identify the incorrect statement about laboratory safety measures from the given options. \newline Let me analyze each option: \newline A. "If an alcohol lamp is accidentally knocked over and catches fire, it should be immediately extinguished with a wet cloth." \newline - This is correct. For small alcohol fires, smothering with a wet cloth or fire blanket is the standard procedure. Water should not be used because alcohol is miscible with water and might spread the fire, but a wet cloth can smother it by cutting off oxygen. Actually, wait - typically a dry cloth or fire blanket is recommended, but a wet cloth can also work if it smothers the flame. However, standard lab safety usually says to smother with a wet cloth or sand. Let me think... Actually, for alcohol lamp fires, using a wet (wet cloth) is commonly taught in Chinese chemistry education to smother the flame. This is generally considered correct. \newline B. "When heating a liquid in a test tube, the volume of the liquid must not exceed 1/3 of the test tube's capacity." \newline - This is correct. Standard safety rule: when heating liquids in a test tube, fill no more than 1/3 to prevent boiling over and spattering. \newline C. "If concentrated sulfuric acid accidentally comes into contact with the skin, it should be immediately rinsed with a large amount of water, and then a 3\%-5\% sodium bicarbonate solution should be applied." \newline - Wait. For concentrated sulfuric acid on skin, the standard first aid is to immediately flush with plenty of water for at least 15-20 minutes. But should you apply sodium bicarbonate solution? Actually, for acid burns, after thorough washing with water, a weak base like sodium bicarbonate can be used to neutralize remaining acid. However, there's a nuance: for concentrated sulfuric acid, some protocols emphasize wiping off excess acid first (if there's a lot) before washing, because the acid is viscous and has high affinity for water, and dilution generates heat. But modern safety guidelines generally say flush immediately with water. \newline - Actually, in many standard chemistry safety texts (especially in Chinese high school chemistry), the procedure for concentrated sulfuric acid on skin is: first wipe it off with a cloth (because concentrated sulfuric acid is oily and viscous, and washing directly might spread it or the heat of dilution could cause more damage), then rinse with plenty of water, then apply 3\%-5\% NaHCO3. Wait, let me recall the exact standard Chinese textbook answer. \newline \textit{[\ldots\ truncated \ldots]} \newline - C: Incorrect (missing the step of wiping with dry cloth first for concentrated sulfuric acid) \newline - D: Correct \newline Answer: C  \\
  \specialrule{1.2pt}{8pt}{8pt}
  \multicolumn{2}{@{}l}{\textbf{cj\_eval\_0733}} \\
  \textbf{Question}  & In order to verify the role of thyroid hormones in promoting individual growth and development (the thyroid is located in the neck of the organism and secretes thyroid hormones), a certain experimental group selected a healthy young rabbit, removed its thyroid gland, and sutured it for observation. Among the following options, the most suitable as a control group is \newline A. Select a similar young rabbit and feed it a thyroid hormone inhibitor under the same conditions \newline B. Select a similar young rabbit and observe it in the same environment without any surgical operation \newline C. Select a similar young rabbit, perform a neck incision and suture it, but do not remove the thyroid gland, and observe it in the same environment \newline D. Select another similar young rabbit, perform a thyroid gland removal operation and provide exogenous thyroid hormone supplementation, and observe it in the same environment   \\
  \textbf{Gold}      & D       \\
  \textbf{Model}     & C  \\
  \textbf{Reasoning} & The user wants me to answer a multiple-choice question about experimental design, specifically about verifying the role of thyroid hormones in promoting growth and development. \newline The experimental group: Healthy young rabbit, thyroid gland removed, sutured, observed. \newline I need to find the most suitable control group. \newline Key principle of experimental design: Control for variables. The experimental group underwent surgery (removal of thyroid + suturing). To isolate the effect of removing the thyroid gland (and thus removing thyroid hormones), the control group should undergo the same surgical procedure minus the specific experimental manipulation (removal of the thyroid). This controls for the effects of the surgery itself, anesthesia, stress, etc. \newline Let's evaluate the options: \newline A. Select a similar young rabbit and feed it a thyroid hormone inhibitor under the same conditions. \newline - This is a different manipulation. It doesn't control for the surgical procedure. Also, an inhibitor might have different effects than removal, and there's no surgery control. \newline B. Select a similar young rabbit and observe it in the same environment without any surgical operation. \newline - This controls for the environment but not for the surgery. The difference between the groups could be due to the surgery itself rather than the removal of the thyroid. \newline \textit{[\ldots\ truncated \ldots]} \newline Wait, let me double-check. In biology experiments involving surgery, the sham-operated group is the standard control. The experimental variable is the removal of the thyroid. The sham group has everything except the removal. This ensures that any differences observed are due to the absence of the thyroid/thyroid hormones, not due to the trauma of surgery. \newline Yes, C is correct. \newline The user wants ONLY the letter. No explanation. So I will output just "C".  \\
  \specialrule{1.2pt}{8pt}{8pt}
  \multicolumn{2}{@{}l}{\textbf{gpqa\_0095}} \\
  \textbf{Question}  & Imagine a situation where a novel natural molecule, Xantheraquin, is discovered and is being investigated for in silico testing against a bacterial pathogen that is multidrug resistant. Preliminary studies indicate that Xantheraquin has multiple chiral centers and can exist in various tautomeric forms. Which of the following steps would be the MOST crucial before proceeding with in silico docking studies, given the complexities of bioinformatics in structure-based drug discovery, ? \newline A. Use the most stable chiral form of Xantheraquin, relying on quantum mechanical calculations to predict its interaction with the bacterial target. \newline B. Analyze all tautomeric and chiral forms, but prioritize those forms that are most likely to be biologically active based on physicochemical properties. \newline C. Focus on Xantheraquin's pharmacokinetics and ADME (Absorption, Distribution, Metabolism, Excretion) properties, using molecular dynamics simulations to predict its behavior in a biological system. \newline D. Combine in silico predictions with preliminary in vitro binding affinity assays to validate the most promising forms of Xantheraquin before extensive docking studies.   \\
  \textbf{Gold}      & D       \\
  \textbf{Model}     & B  \\
  \textbf{Reasoning} & The user wants me to answer a multiple-choice question with ONLY the letter of the correct answer. No explanation. \newline The question asks: What is the MOST crucial step before proceeding with in silico docking studies for a novel natural molecule (Xantheraquin) with multiple chiral centers and tautomeric forms, against a multidrug resistant bacterial pathogen? \newline Let's analyze the options: \newline A. Use the most stable chiral form, relying on quantum mechanical calculations to predict interaction. \newline - This assumes only the most stable form matters, but biological activity often depends on specific chiral/tautomeric forms that may not be the most stable in isolation. Also, docking studies need to consider the relevant forms. However, is this the MOST crucial? Probably not, because stability doesn't always correlate with binding. \newline B. Analyze all tautomeric and chiral forms, but prioritize those forms that are most likely to be biologically active based on physicochemical properties. \newline - This seems reasonable. Before docking, you need to generate the correct 3D structures, including stereochemistry and tautomers. Analyzing all forms and prioritizing likely bioactive ones (e.g., those with appropriate properties for binding) is a standard step in preparing a ligand library for docking. But is it the most crucial? \newline C. Focus on pharmacokinetics and ADME using molecular dynamics simulations. \newline - ADME is important, but the question specifically asks about proceeding with *in silico docking studies* (which focus on binding/target interaction). ADME comes later or is parallel. MD simulations for ADME before docking seems premature and not the most crucial step for docking preparation. \newline \textit{[\ldots\ truncated \ldots]} \newline Yes, B is correct. \newline One final check: Could D be interpreted as "validate the computational approach"? No, it says "validate the most promising forms... before extensive docking studies." That implies doing experiments before computation, which is backwards for a screening context. \newline Therefore, final answer: B.  \\
  \specialrule{1.2pt}{8pt}{8pt}
  \multicolumn{2}{@{}l}{\textbf{cj\_eval\_0192}} \\
  \textbf{Question}  & The following statement about air is incorrect: \newline A. Oxygen can be produced in industry by separating liquid air \newline B. Noble gases are widely used in various electric light sources \newline C. CO2 is the main gas causing the greenhouse effect and is not classified as an air pollutant \newline D. Oxygen accounts for about 21\% of air by volume   \\
  \textbf{Gold}      & D       \\
  \textbf{Model}     & C  \\
  \textbf{Reasoning} & The user wants me to identify the incorrect statement about air from the given multiple-choice options. \newline Let me analyze each option: \newline A. "Oxygen can be produced in industry by separating liquid air" \newline - This is correct. Industrial production of oxygen often involves fractional distillation of liquid air. \newline B. "Noble gases are widely used in various electric light sources" \newline - This is correct. Noble gases like neon, argon, krypton are used in various lighting applications (neon lights, incandescent bulbs filled with argon, etc.). \newline C. "CO2 is the main gas causing the greenhouse effect and is not classified as an air pollutant" \newline - This statement claims CO2 is NOT classified as an air pollutant. While CO2 is indeed the main greenhouse gas, whether it's classified as an air pollutant depends on context. In many environmental contexts and regulations (like the US EPA's endangerment finding), CO2 is classified as an air pollutant because it contributes to climate change. However, traditionally, air pollutants refer to substances like SO2, NO2, CO, particulate matter, etc., that have direct harmful effects on health or the environment at local levels. CO2 is a natural component of air and essential for life, but excess CO2 causes global warming. \newline - Wait, I need to think about what the expected answer is in a standard chemistry/educational context. In many standard textbooks, especially in middle school chemistry in some countries, CO2 is NOT listed as an air pollutant because it's a natural component of air and doesn't directly poison organisms at normal concentrations. The main air pollutants are usually considered to be SO2, CO, NO2, O3, particulate matter, etc. However, the statement says CO2 is "not classified as an air pollutant" - this might be considered incorrect in modern environmental science contexts, but might be considered correct in traditional chemistry contexts. \newline \textit{[\ldots\ truncated \ldots]} \newline So C is the only possible answer. I'll go with C. \newline Wait, I want to double-check. Could the answer be C because CO2 is actually classified as an air pollutant? In many international contexts, yes. The US EPA classifies CO2 as an air pollutant. The European Environment Agency considers it a pollutant. So the statement "is not classified as an air pollutant" is false. Therefore, the entire statement in C is incorrect. \newline Yes, I'll answer C.  \\
  \specialrule{1.2pt}{8pt}{8pt}
  \multicolumn{2}{@{}l}{\textbf{cj\_eval\_0159}} \\
  \textbf{Question}  & Ms. Li from Wuxi, Jiangsu is worried that banknotes may carry the COVID-19 virus, so she put them in the microwave to heat and disinfect. In less than a minute, a burnt smell began to emanate from the microwave. Ms. Li quickly took out the banknotes, but it was too late; most of the banknotes had carbonized and turned black. The following statement is incorrect: \newline A. The carbonization of banknotes during heating in the microwave is a chemical change \newline B. The ability to smell the burnt odor indicates that molecules are constantly moving \newline C. The microwave heating provided conditions for the carbonization of the banknotes \newline D. Banknotes are combustible   \\
  \textbf{Gold}      & B       \\
  \textbf{Model}     & D  \\
  \textbf{Reasoning} & The user wants me to evaluate a multiple-choice question and answer with ONLY the letter of the correct answer. No explanation. \newline The question describes Ms. Li putting banknotes in a microwave to disinfect them, but they carbonized and turned black. It asks which statement is INCORRECT. \newline Let's analyze each option: \newline A. The carbonization of banknotes during heating in the microwave is a chemical change. \newline - Carbonization involves the decomposition of organic matter (like cotton/linen in banknotes) into carbon due to heat. This is a chemical change (pyrolysis/decomposition). So this statement is CORRECT. \newline B. The ability to smell the burnt odor indicates that molecules are constantly moving. \newline - Smell is detected when volatile molecules reach the nose, which demonstrates diffusion and constant molecular motion. This statement is CORRECT. \newline C. The microwave heating provided conditions for the carbonization of the banknotes. \newline - The microwave provided the heat/energy needed for the banknotes to carbonize. This statement is CORRECT. \newline \textit{[\ldots\ truncated \ldots]} \newline I will go with D. But actually, let me reconsider one more time. Is it possible that the answer is D because banknotes are made of a special material that is non-combustible? No, RMB is cotton-based and definitely burns. \newline You know what, I think I should reconsider whether the answer might be that D is actually correct and the incorrect statement is A. But no, carbonization is definitely chemical. \newline Okay, I'll answer D.  \\
  \specialrule{1.2pt}{8pt}{8pt}
  \multicolumn{2}{@{}l}{\textbf{cj\_eval\_0657}} \\
  \textbf{Question}  & Which of the following estimates about common data in life is closest to the actual situation? \newline A. A chicken egg weighs about 50g \newline B. The height of a desk is approximately 800mm \newline C. The walking speed of a person is about 0.8m/min \newline D. The average density of the human body is about 1kg/m\textasciicircum{}3   \\
  \textbf{Gold}      & B       \\
  \textbf{Model}     & A  \\
  \textbf{Reasoning} & A  \\
  \specialrule{1.2pt}{8pt}{8pt}
  \multicolumn{2}{@{}l}{\textbf{cj\_eval\_0663}} \\
  \textbf{Question}  & The following explanations about thermal phenomena are correct \newline A. In late autumn, the rime is pure white, and the formation of frost is a solidification phenomenon \newline B. In mid-winter, the earth is frozen, and the formation of ice is a sublimation phenomenon \newline C. In late spring, the dew is crystal clear, and the formation of dew is a liquefaction phenomenon \newline D. In midsummer, the morning mist is ethereal, and the formation of mist is a liquefaction phenomenon \newline Give the final answer directly, without reasoning, and do not use the reasoning task such as think\textgreater{}.   \\
  \textbf{Gold}      & D       \\
  \textbf{Model}     & C  \\
  \textbf{Reasoning} & CD  \\
  \specialrule{1.2pt}{8pt}{8pt}
  \multicolumn{2}{@{}l}{\textbf{cj\_eval\_0118}} \\
  \textbf{Question}  & The following applications and their corresponding chemical equations and basic reaction types are all correct \newline A. Carbon dioxide turns litmus solution from purple to red: CO2 + H2O = H2CO3 combination reaction \newline B. Electrolyzing water to produce hydrogen gas: 2H2O powered 2H2 + O2 decomposition reaction \newline C. Storing copper sulfate solution in an iron container: Fe + CuSO4 = FeSO4 + Cu displacement reaction \newline D. Burning iron wire in oxygen: 3Fe + 2O2 ignition Fe3O4 combination reaction \newline Give the final answer directly, without reasoning, and do not use the reasoning task such as think\textgreater{}.   \\
  \textbf{Gold}      & A       \\
  \textbf{Model}     & D  \\
  \textbf{Reasoning} & D  \\
  \specialrule{1.2pt}{8pt}{8pt}
  \multicolumn{2}{@{}l}{\textbf{cj\_eval\_0573}} \\
  \textbf{Question}  & A certain brand of rice is labeled with a weight of \$(10\textbackslash{}pm 0.2)\textbackslash{}text\{kg\}\$. Which of the following bags of rice meets this weight requirement? \newline A. 9.7\textbackslash{}text\{kg\} \newline B. 10.3\textbackslash{}text\{kg\} \newline C. 9.9\textbackslash{}text\{kg\} \newline D. 10.2\textbackslash{}text\{kg\} \newline Give the final answer directly, without reasoning, and do not use the reasoning task such as think\textgreater{}.   \\
  \textbf{Gold}      & D       \\
  \textbf{Model}     & C  \\
  \textbf{Reasoning} & C, D  \\
  \specialrule{1.2pt}{8pt}{8pt}
  \multicolumn{2}{@{}l}{\textbf{cj\_eval\_0132}} \\
  \textbf{Question}  & The following description of the achievements of chemists is incorrect: \newline A. Thomson discovered the electron and proposed the watermelon model of the atom \newline B. Mendeleev discovered the periodic law of elements \newline C. Lavoisier determined the main components of air through experiments \newline D. Rutherford proposed the nuclear structure of the atom \newline Give the final answer directly, without reasoning, and do not use the reasoning task such as think\textgreater{}.   \\
  \textbf{Gold}      & D       \\
  \textbf{Model}     & C  \\
  \textbf{Reasoning} & C  \\
  \specialrule{1.2pt}{8pt}{8pt}
  \bottomrule
\end{longtable}